\documentclass[11pt,a4paper]{article}
\usepackage[utf8]{inputenc}
\usepackage[T1]{fontenc}
\usepackage[french]{babel}
\usepackage{geometry}
\usepackage{helvet}
\usepackage{titlesec}
\usepackage{parskip}
\usepackage{xcolor}
\usepackage{hyperref}

\geometry{a4paper, margin=2.5cm}
\renewcommand{\familydefault}{\sfdefault}

\definecolor{primary}{RGB}{13, 27, 62}
\definecolor{secondary}{RGB}{65, 139, 255}

\titleformat{\section}{\color{primary}\Large\bfseries}{\thesection}{1em}{}
\titleformat{\subsection}{\color{secondary}\large\bfseries}{\thesubsection}{1em}{}

\hypersetup{
    colorlinks=true,
    linkcolor=primary,
    filecolor=magenta,      
    urlcolor=secondary,
    pdftitle={Cahier des Charges - Conversion Excel vers Word},
}

\begin{document}

\begin{titlepage}
    \centering
    \vspace*{4cm}
    {\Huge \textbf{\textcolor{primary}{Cahier des Charges Opérationnel}}} \\[1cm]
    {\LARGE Application de Conversion Sécurisée Excel vers Word} \\[2cm]
    {\large \textbf{Version 1.1 -- Projet Interne}} \\[3cm]
    
    \vfill
    {\large Août 2026}
\end{titlepage}

\tableofcontents
\newpage

\section{Contexte et Problématique}
Au sein de l'environnement opérationnel de l'entreprise, les équipes métiers manipulent quotidiennement un volume important de données structurées, consolidées et traitées sous format Microsoft Excel. Ces données, allant des reportings en ressources humaines aux états financiers, doivent régulièrement être communiquées à des instances directionnelles ou des partenaires extérieurs sous une forme narrative, formelle et rédigée, généralement au format Microsoft Word.

Le processus actuel repose entièrement sur des interventions manuelles chronophages. La transformation de l'information exige qu'un collaborateur ouvre un fichier Excel complexe, en filtre le contenu pour ensuite le transférer par des opérations de copie et de collage répétitives au sein d'un modèle prédéfini. Cette approche soulève un défi majeur : comment automatiser la conversion et le formatage d'un jeu de données Excel vers un document Word formel, tout en garantissant la sécurité, la traçabilité des données d'entreprise et la conformité de la mise en forme.

\section{Étude Comparative de l'Existant}
Le marché propose diverses solutions de conversion, cependant celles-ci demeurent profondément inadaptées à un contexte d'entreprise soumis à de fortes contraintes sécuritaires. 

Les convertisseurs publics accessibles en ligne constituent une faille de sécurité inacceptable. En téléversant des documents sur de tels espaces, l'entreprise perd la souveraineté sur ses données sensibles, compromettant le respect de la confidentialité et des réglementations en vigueur. De plus, le résultat produit relève davantage d'une impression statique imposée que d'un document Word structuré et véritablement éditable selon un formatage rigoureux.

Les alternatives reposant sur des macros développées localement, ou des outils d'automatisation bureautique sur le poste de l'utilisateur, pallient le risque de fuite des données mais engendrent des problématiques de maintenance considérables. La pérennité de ces scripts est mise en péril à chaque mise à jour logicielle. En outre, ce modèle d'exécution décentralisée empêche toute forme d'audit et de traçabilité des opérations à l'échelle du système d'information.

Bien que formidables en termes de connectivité, les offres d'automatisation cloud professionnelles requièrent souvent des licences onéreuses et font transiter les données au travers de services externes, sans oublier la complexité induite pour l'utilisateur final qui souhaite simplement obtenir un document exploitable à partir de son fichier.

\section{Proposition Fonctionnelle et Valeur Ajoutée}
En réponse aux limites constatées, il est proposé de concevoir une application de conversion sur-mesure, déployée en interne. Cette approche internalisée garantit par conception une souveraineté totale, les fichiers traités par les collaborateurs subsistant à l'intérieur du périmètre réseau sécurisé de l'entreprise.

Une dimension cruciale de la proposition réside dans la gestion des accès et de la traçabilité. Chaque opération, soumise à l'authentification préalable des utilisateurs, générera une empreinte inaltérable permettant à l'administration de s'assurer de l'identité du demandeur, de l'horodatage et de la nature de la transaction. L'interface applicative devra par ailleurs se distinguer par une grande simplicité, facilitant une adhésion immédiate des équipes grâce à un mécanisme de téléversement instinctif.

Sur le plan du traitement des fichiers, la solution logicielle visera à analyser la structure des classeurs Excel pour reconstituer intelligemment l'arborescence des tableaux et sections appropriées au format Word natif, aboutissant à un rendu soigné et directement prêt à la diffusion professionnelle.

\section{Décision sur l'Architecture et la Stack Technique}
Le pilotage de la solution applicative sera assuré par une architecture séparant distinctement le client et le serveur. L'étude de la plateforme technique optimale s'est naturellement orientée vers un écosystème éprouvé en milieu industriel : le couplage entre l'environnement Java et une interface utilisateur de nouvelle génération.

Le traitement intensif des fichiers nécessite un socle robuste capable d'analyser le code des documents sans dépendre des suites logicielles classiques. Le framework Spring Boot a été sélectionné pour sa maturité incontestable dans la gestion des API sécurisées en entreprise, et pour son intégration sans faille avec le module Apache POI. Cette bibliothèque Java de référence apporte toutes les garanties de performance requises pour l'instanciation de documents riches au format de la bureautique moderne. 

L'expérience applicative côté client nécessitera, quant à elle, l'intervention du cadriciel React associé à l'empaqueteur Vite. Ensemble, ils autorisent le développement d'interfaces fluides et robustes, garantes d'une transition en douceur pour les opérationnels, tout en facilitant l'intégration des barrières de sécurité indispensables telles que la transmission des jetons cryptographiques de courte durée pour la connexion.

L'infrastructure de données s'appuiera sur PostgreSQL pour l'intégrité de son suivi des interactions, assurant le bon ordonnancement sans risque d'ingérer l'intégralité des fichiers, ces derniers demeurant isolés temporellement sur un espace de stockage volatil du backend en attente de restitution au demandeur. En adoptant ce schéma technologique, la solution finalisée alliera vélocité applicative et conformité institutionnelle.

\end{document}
